\section{The optimal cycle}\label{sec:cycle}

The results of Sections~1--8 provide the three stages of a closed algebraic loop:
\[
M = (\Rpos, \times)
\xrightarrow{\;\log\;}
A = (\R, +)
\xrightarrow{\;\text{Fourier}\;}
\widehat{A} = (\Z, +)
\xrightarrow{\;\exp\;}
M.
\]
We now formalize this loop and its obstruction to closure.

\begin{definition}[The additive-multiplicative cycle]\label{def:cycle}
Let $f \in L^1(\bT)$ with $f > 0$ a.e.\ and $\log f \in L^1(\bT)$. Write
\[
c_k = \int_0^1 (\log f(\theta))\, e^{-2\pi i k\theta}\, d\theta, \qquad k \in \Z,
\]
for the Fourier coefficients of $\log f$. The \emph{additive-multiplicative cycle} of $f$ consists of:
\begin{enumerate}[label=(\roman*)]
\item The \emph{DC component}: $c_0 = \int_0^1 \log f(\theta)\, d\theta$, yielding the geometric mean $G(f) = e^{c_0}$.
\item The \emph{AC component}: $\{c_k\}_{k \neq 0}$, yielding the \emph{holonomy}
\begin{equation}\label{eq:holonomy}
H(f) = \sum_{k=1}^{\infty} k\, |c_k|^2 \in [0, \infty].
\end{equation}
\item The \emph{Szeg\H{o} constant}: $E(f) = \exp(H(f)) \in [1, \infty]$.
\end{enumerate}
The holonomy $H(f)$ is the one-sided $H^{1/2}(\bT)$ Sobolev seminorm of $\log f$ (summing over positive frequencies; for real-valued $\log f$, this equals $\frac{1}{2}\|\log f\|_{H^{1/2}(\bT)}^2$) and measures the cost of traversing the cycle.
\end{definition}

\begin{theorem}[Holonomy trichotomy]\label{thm:holonomy-trichotomy}
Let $f$ be as in \cref{def:cycle}. Then exactly one of the following holds:
\begin{enumerate}[label=(\roman*)]
\item $H(f) = 0$, equivalently $E(f) = 1$, if and only if $c_k = 0$ for all $k \geq 1$, if and only if $f$ is constant a.e.\ (the Gaussian spectral density). The cycle closes perfectly.
\item $0 < H(f) < \infty$, equivalently $1 < E(f) < \infty$: the Strong Szeg\H{o} class. The cycle is non-trivial with finite cost, and $D_n(f) \sim G(f)^n \cdot E(f)$.
\item $H(f) = \infty$, equivalently $E(f) = \infty$: the cycle breaks. The ratio $D_n(f) / G(f)^n$ diverges or oscillates without a finite limit.
\end{enumerate}
\end{theorem}

\begin{proof}
Part~(i) is immediate from the definition: $H(f) = 0$ requires $c_k = 0$ for all $k \geq 1$, so $\log f = c_0$ is constant, hence $f = e^{c_0}$ is constant. Parts~(ii) and~(iii) are the content of the Strong Szeg\H{o} Limit Theorem (\cref{thm:szego} in \S\ref{sec:hardy}).
\end{proof}

\begin{remark}[The Gaussian as unique fixed point]
The Gaussian spectral density $f \equiv \sigma^2$ is simultaneously:
\begin{enumerate}[label=(\alph*)]
\item the Fourier eigenfunction (\cref{thm:gaussian-fixed}),
\item the heat kernel (\cref{thm:heat-semigroup}),
\item the unique outer function seed ($c_k = 0$ for $k \neq 0$),
\item the unique spectral density with $E(f) = 1$ (\cref{thm:holonomy-trichotomy}(i)).
\end{enumerate}
These four characterizations are manifestations of a single fact: the Gaussian is the fixed point of the optimal cycle.
\end{remark}

\begin{definition}[Constrained optimal cycle]\label{def:optimal-cycle}
Let $\Omega$ denote the space of spectral densities $f \in L^1(\bT)$ with $f > 0$ a.e.\ and $\log f \in L^1(\bT)$. For each $t \geq 0$, define the \emph{feasible set at lag $t$}:
\[
\Omega_t = \bigl\{f \in \Omega : \widehat{f}(k) = r_k \text{ for } |k| \leq t\bigr\},
\]
where $r_0, r_1, \ldots$ are given autocovariance values (observed data). The \emph{constrained holonomy} is
\[
H_t^* = \inf_{f \in \Omega_t} H(f).
\]
The minimizer $f_t^*$, when it exists, is the spectral density of minimum holonomy compatible with the observations.
\end{definition}

\begin{proposition}[Monotonicity of the constrained cycle]\label{prop:holonomy-monotone}
In the notation of \cref{def:optimal-cycle}:
\begin{enumerate}[label=(\roman*)]
\item \textbf{Shrinking feasible sets}: $\Omega_0 \supseteq \Omega_1 \supseteq \Omega_2 \supseteq \cdots$.
\item \textbf{Monotone holonomy}: $H_0^* \leq H_1^* \leq \cdots \leq H(f_{\mathrm{true}})$, with equality in the limit if $f_{\mathrm{true}} \in \Omega$.
\item \textbf{Maximum-entropy minimizer}: the minimizer $f_t^*$ is the maximum-entropy spectral density compatible with $r_0, \ldots, r_t$, i.e., the autoregressive AR$(t)$ spectral density (Burg's theorem \cite{Burg}).
\item \textbf{Dual convergence}: the prediction error $\sigma_t^2 = \exp\bigl(\int_{\bT} \log f_t^*\bigr)$ decreases monotonically to $\sigma_{\mathrm{true}}^2$, while $H_t^*$ increases monotonically to $H(f_{\mathrm{true}})$---convergence from opposite directions.
\end{enumerate}
\end{proposition}

\begin{proof}
Part~(i): each additional constraint $\widehat{f}(t+1) = r_{t+1}$ restricts the feasible set.

Part~(ii): immediate from~(i), since the infimum over a smaller set is at least as large.

Part~(iii): by the Szeg\H{o}--Verblunsky identity (\cref{prop:verblunsky}), the Verblunsky coefficients $\alpha_0, \ldots, \alpha_{t-1}$ are uniquely determined by the autocovariances $r_0, \ldots, r_t$ via the Levinson recursion. They are the same for all $f \in \Omega_t$. The AR$(t)$ model is characterized by $\alpha_k = 0$ for all $k \geq t$. Split $H(f)$ using \cref{prop:verblunsky}(iii):
\[
H(f) = \underbrace{-\sum_{k=1}^{t-1} k\log(1-|\alpha_k|^2)}_{\text{constant on }\Omega_t} \;+\; \underbrace{\left(-\sum_{k=t}^{\infty} k\log(1-|\alpha_k|^2)\right)}_{\geq\, 0,\;\; =\, 0 \;\text{iff}\; \alpha_k = 0\;\forall\, k \geq t}.
\]
Similarly, for the geometric mean (\cref{prop:verblunsky}(ii)):
\[
\log G(f) = \underbrace{\log r_0 + \sum_{k=0}^{t-1}\log(1-|\alpha_k|^2)}_{\text{constant on }\Omega_t} \;+\; \underbrace{\sum_{k=t}^{\infty}\log(1-|\alpha_k|^2)}_{\leq\, 0,\;\; =\, 0 \;\text{iff}\; \alpha_k = 0\;\forall\, k \geq t}.
\]
Since $-\log(1-|\alpha_k|^2) \geq 0$ with equality iff $\alpha_k = 0$: $H(f) \geq H(f_{\mathrm{AR}(t)})$ and $G(f) \leq G(f_{\mathrm{AR}(t)})$, with equality in both iff $\alpha_k = 0$ for all $k \geq t$, i.e., $f = f_{\mathrm{AR}(t)}$. This is Burg's theorem \cite{Burg}: the AR$(t)$ density simultaneously minimizes holonomy and maximizes the geometric mean (entropy) on $\Omega_t$.

Part~(iv): since $f_t^* = f_{\mathrm{AR}(t)}$, we have $\sigma_t^2 = G(f_t^*)$ and $H_t^* = H(f_t^*)$. As $t$ increases: each additional Verblunsky coefficient $\alpha_t \neq 0$ (generically) adds a positive term to $H$ and a negative term to $\log G$, so $H_t^*$ increases and $\sigma_t^2$ decreases monotonically. These are dual aspects of the same convergence: the Verblunsky coefficients apportion the true spectral density's holonomy into constrained (known) and free (unknown) parts, and each observation converts one coefficient from free to constrained.
\end{proof}

\begin{remark}[Connection to GPI]
The filtration $\mathcal{G}_t = \sigma(r_0, \ldots, r_t)$ generated by observed autocovariances is a functor in $\Proj_Y$ (\cref{def:gpi-category}). The constrained holonomy $H_t^*$ is the evaluation of this functor at the holonomy functional, and its monotonicity (\cref{prop:holonomy-monotone}(ii)) is an instance of the residual monotonicity of \cref{cor:monotonicity}. The dual convergence---prediction error decreasing, holonomy increasing---is the content of the in-context GPI theorem (\cref{thm:in-context-gpi}) applied to the spectral density space.
\end{remark}

\begin{remark}[Measurement, estimation, and prediction]\label{rem:three-layers}
The re-cycle operates through three layers, each formalized by the preceding results.

\emph{Layer~1: Measurement.}
Build the $\sigma$-algebra $\mathcal{G}_t = \sigma(r_0, \ldots, r_t)$ from observed autocovariances. Without measurement, $\mathcal{G} = \{\varnothing, \Omega\}$, and by \cref{prop:cond-exp-optimal}:
\[
\mathbb{E}[Y \mid \mathcal{G}] = \mathbb{E}[Y]
\]
---the global mean, maximum ignorance. In cycle language: the feasible set $\Omega_0$ is unconstrained, so $H_0^* = 0$ trivially---not because the true density is Gaussian, but because no data constrains the minimizer away from the constant density. The prediction is pure DC: $G(f)$ without $E(f)$.

\emph{Layer~2: Estimation (re-cycle).}
Each measurement refines $\mathcal{G}_t$. By \cref{cor:monotonicity}, the residual $R_Y(\mathcal{G}_t)$ decreases monotonically. In spectral language: $f_t^*$ (the AR$(t)$ minimizer of \cref{prop:holonomy-monotone}) sharpens, $\nu$ increases, $t_\nu$ thins. This is slow but guaranteed: $O(1/\sqrt{n})$ by the CLT rate.

\emph{Layer~3: Prediction (direct computation).}
Given the estimated spectral density $f_t^*$, compute consequences exactly---the forward problem on known structure. This layer is only as good as Layer~2: prediction inherits the estimation uncertainty, but the forward computation is exact, so input uncertainty maps to output uncertainty in a fully characterized way.

The closed operational loop
\begin{gather*}
\text{Measure}
\xrightarrow{\;\text{re-cycle}\;}
\text{Estimate}
\xrightarrow{\;\text{direct computation}\;}
\text{Predict} \\
\xrightarrow{\;\text{target next observation}\;}
\text{Measure better}
\end{gather*}
is the optimal cycle applied operationally. Each revolution refines Layer~2: $\nu$ increases, $H$ decreases.
\end{remark}

\begin{example}[The Student $t$-distribution as canonical re-cycle]\label{ex:student-t}
\emph{Student $t$ with $\nu$ degrees of freedom is the re-cycle snapshot at observation $\nu$.}

The Student $t_\nu$ family parameterizes the entire re-cycle trajectory with a single real parameter $\nu \in [1, \infty]$:
{\small
\[
\renewcommand{\arraystretch}{1.4}
\begin{array}{cclcl}
\hline
\nu & \text{Distribution} & H(f_{t_\nu}) & E(f_{t_\nu}) & \text{Cycle state} \\
\hline
\nu = 1 & \text{Cauchy} & \infty & \infty & \text{Broken} \\
1 < \nu \leq 2 & \text{Heavy-tailed } t & \infty & \infty & E[X^2] = \infty \\
2 < \nu < \infty & \text{Moderate } t & \text{finite}, > 0 & > 1 & \text{Non-trivial---learning} \\
\nu \to \infty & \text{Gaussian} & 0 & 1 & \text{Fixed point} \\
\hline
\end{array}
\]
}\end{example}

\begin{proposition}[Student $t$ convergence along the re-cycle]\label{prop:student-convergence}
For $\nu > 0$, define the \emph{spectral density associated to the Student $t_\nu$ distribution} as the periodized autocovariance:
\[
f_{t_\nu}(\theta) = \sum_{k \in \Z} r_k\, e^{2\pi ik\theta}, \qquad r_k = \varphi_{t_\nu}(2\pi k),
\]
where $\varphi_{t_\nu}(\xi) = \frac{2^{1-\nu/2}}{\Gamma(\nu/2)}|\sqrt{\nu}\,\xi|^{\nu/2}K_{\nu/2}(\sqrt{\nu}\,|\xi|)$ is the characteristic function of the $t_\nu$ distribution and $K_{\nu/2}$ is the modified Bessel function of the second kind. For $\nu > 2$, the second moment $E[X^2] = \nu/(\nu-2)$ is finite, $r_k$ decays exponentially in $|k|$, and $f_{t_\nu} > 0$ with $\log f_{t_\nu} \in L^1(\bT)$.
\begin{enumerate}[label=(\roman*)]
\item The critical value is $\nu^* = 2$: the Szeg\H{o} constant satisfies $E(f_{t_\nu}) < \infty$ if and only if $\nu > 2$, corresponding to the threshold where the second moment $E[X^2]$ becomes finite and the autocovariance is well-defined and summable.
\item For $\nu > 2$: $H(f_{t_\nu}) = O(\log\nu / \nu)$ as $\nu \to \infty$. The convergence rate to the Gaussian fixed point is algebraic.
\item The unique minimizer of $H$ over the family $\{f_{t_\nu} : \nu > 2\}$ is $\nu = \infty$ (Gaussian), confirming the cycle's unique fixed point.
\end{enumerate}
\end{proposition}

\begin{proof}[Proof sketch]
(i) For $\nu \leq 2$, the second moment diverges ($E[X^2] = \nu/(\nu-2) \to \infty$ as $\nu \to 2^+$), so $r_0 = \infty$ and the spectral density is not in $L^1$. For $\nu > 2$, the exponential decay of $K_{\nu/2}$ ensures $r_k = O(e^{-c|k|})$ for some $c > 0$, so $f_{t_\nu}$ is smooth and strictly positive, giving $\log f_{t_\nu} \in H^{1/2}(\bT)$ and $E(f_{t_\nu}) < \infty$.

(ii) By the Szeg\H{o}--Verblunsky identity (\cref{prop:verblunsky}), $H(f_{t_\nu}) = -\sum_{k=1}^{\infty}k\log(1-|\alpha_k(\nu)|^2)$. For large $\nu$, the $t_\nu$ characteristic function satisfies $\varphi_{t_\nu}(\xi) = e^{-\xi^2/2}(1 + O(\xi^4/\nu))$, so the spectral density approaches the Gaussian at rate $O(1/\nu)$ uniformly. The Verblunsky coefficients inherit this: $|\alpha_k(\nu)|^2 = O(k^{-2}\nu^{-1})$, giving $H(f_{t_\nu}) \approx \sum_{k=1}^{N} k|\alpha_k|^2 = O(\nu^{-1})\sum_{k=1}^{N} k^{-1} = O(\log\nu/\nu)$ where $N \sim \nu$ is the effective cutoff beyond which $|\alpha_k|$ is negligible.

(iii) By \cref{prop:verblunsky}(iii), $H(f_{t_\nu}) > 0$ for any $\nu < \infty$ (since some $\alpha_k \neq 0$), and $H(f_{t_\infty}) = 0$ (all $\alpha_k = 0$, Gaussian).
\end{proof}

\begin{remark}[Canonicality]
The Student $t$ family is canonical because it parameterizes the \emph{entire} re-cycle trajectory with a single real parameter $\nu \in [1, \infty]$. Every other natural family (tempered stable, variance-gamma, generalized hyperbolic) traverses only part of the trichotomy. The $t$-family is the geodesic through holonomy space.
\end{remark}

\begin{remark}[Only Gauss sees the Gaussian]\label{rem:only-gauss}
This remark has three layers of meaning.

\emph{Historical.} Gauss understood this object; others used it without seeing it. The name ``normal distribution''---never Gauss's term---encodes the false belief that the Gaussian is the \emph{common case}. It is not. It is the \emph{limit}.

\emph{Mathematical.} $E(f) = 1$ if and only if $f$ is Gaussian (\cref{thm:holonomy-trichotomy}). Only the Gaussian system has zero holonomy. Only the Gaussian sees itself perfectly through the cycle. Every other system incurs loss.

\emph{Epistemological.} $E(f_{t_\nu}) = 1$ requires $\nu = \infty$: the observer must \emph{become} the system. At $\nu = \infty$ you are no longer observing---you are the cycle itself. The map has become the territory.

The CLT does not say ``things converge to normal.'' It says: \emph{the re-cycle converges to the cycle}. The Student $t_\nu$ is the re-cycle at finite time. The Gaussian is the cycle itself. Every actual measurement is $t_\nu$ with some finite $\nu$. The Gaussian is what they approach but none of them are---the $c$ of probability theory: unreachable, structural, defining everything by being the limit nothing reaches.
\end{remark}

